\documentclass[10pt,a4paper]{amsart}
\usepackage[margin=19mm]{geometry}
\usepackage{amsmath,amssymb,mathtools}
\usepackage[scr=rsfs,cal=euler]{mathalfa}
\usepackage{xfrac}
\usepackage[unicode,hidelinks,bookmarks=false]{hyperref}
\newcommand{\ie}{i.e.\ }
\newcommand{\cf}{cf.\ }
\newcommand{\resp}{respectively\ }
\newcommand{\Subsection}{\subsection}
\providecommand{\nobreakdash}{\nobreak\hbox{-}}
\setlength{\emergencystretch}{2em}
\multlinegap0pt
\usepackage[shortlabels]{enumitem}
\usepackage{mathtools}
\usepackage{amssymb}
\usepackage{float}
\usepackage{csquotes}
\newtheorem{lemma}{Lemma}
\newtheorem{claim}{Claim}
\newtheorem{subclaim}{Claim}
\title{Quagmires and large Suslin forests}
\author{Lorenzo Notaro}
\date{Source: arXiv:2608.28505v1 (CC BY 4.0)}
\begin{document}
\maketitle

\noindent\emph{Source and licence.} Quagmires and large Suslin forests. Version arXiv:2608.28505v1 (CC BY 4.0), distributed by arXiv under the Creative Commons Attribution 4.0 International licence, http://creativecommons.org/licenses/by/4.0/, as recorded in the arXiv licence metadata for this exact version and shown on the abstract page https://arxiv.org/abs/2608.28505v1. Full licence notice: \texttt{licenses/arxiv-2608.28505-CC-BY-4.0.txt}.\par\smallskip

\noindent\textit{Editorial scope.} Counted unit: the article's claim claim:5 (source0.tex lines 521--527). The article's complete original proof of it is delivered with it (source0.tex lines 528--592, one \texttt{proof} environment of 65 lines). No mathematics is written, repaired or re-derived: the statement and the proof are the article's own lines, in the article's own notation, and no retained line is substituted or re-flowed.  Retained uncounted as the source-local context the proof needs: lines 271--279.  Nothing was omitted from any retained span, so the package carries no omission record.  No retained line contains a citation, so the wrapper carries no bibliography and no bibliography entry.  No retained line contains a figure environment, an image-inclusion command or drawing code, so no image is delivered and the package declares no asset dependency.  Editorial formatting only, in addition: an AMS \texttt{amsart} preamble that restates the article's own declarations and macros used by the retained lines, namely the environments \texttt{lemma}, \texttt{claim}, \texttt{subclaim} on separate counters, together with the standard packages those lines need.  The retained environments print in this document as Lemma 1, Claim 1, Subclaim 1 in order of appearance, whereas the article prints the same items with the numbering of its own full document.  The complete native source of the article as distributed in the arXiv e-print for this exact version is bundled unchanged as \texttt{sources/208793/source0.tex}.
\begin{lemma}\label{lemma:basicquag}
Let $(T, <, \lhd, Q)$ be a $(\kappa, S)$-quagmire. Then the following hold:
\begin{enumerate}[label={\upshape (\arabic*)}]
\item $|T_\kappa| = \kappa^+$ and every element of $[T_{\kappa}]^{\le \kappa}$ has an $\lhd$-upper bound;
\item for every $x \in T$ and every $\alpha < \mathrm{ht}(x)$, the map $Q(\cdot, x) \upharpoonright T_\alpha$ is an order-embedding;
\item for every $(x, y) \in \mathrm{dom}(Q)$ and $\alpha$ with $\mathrm{ht}(x) < \alpha < \mathrm{ht}(y)$, $\pi_\alpha(Q(x, y)) = Q(x, \pi_\alpha(y))$;
\item for every $x, y$ and $\alpha < \beta < \mathrm{ht}(x)$, if $x = Q(\pi_\alpha(x), y)$, then $x = Q(\pi_\beta(x), y)$.
\end{enumerate}
\end{lemma}

\begin{claim}\label{claim:5}
Let $\alpha \in S'$ and $A \in V_\alpha$. If $\alpha$ satisfies $(\dagger)$ with witnesses $N, Y$, then:
\begin{enumerate}[label={\upshape (\arabic*)}]
\item $Y \cap \pi_{\alpha \lambda}^{-1}(A) \cap N \in N$,
\item $A \cap X^\alpha_0 \in V_\alpha$.
\end{enumerate}
\end{claim}

\begin{proof}

Suppose that $\alpha \in S'$ satisfies $(\dagger)$ with witnesses $N, Y$, and fix $A \in V_\alpha$, towards showing that (1) and (2) hold. Since both $V_\alpha$ and $N$ are closed under ${<}\kappa$-unions, we can assume without loss of generality that $A = D(x, \beta, \xi)$ for some $x \in T_\alpha$ and $\beta, \xi < \alpha$. 

So fix $\beta, \xi < \alpha$ and $x \in T_\alpha$ and  let 
\[
K \coloneqq Y \cap \pi_{\alpha \lambda}^{-1}(D(x, \beta, \xi)) \cap N.
\] 

By Lemma~\ref{lemma:basicquag}, we know that $Y$ has a $\lhd$-upper bound. By elementarity of $N$, we can find such an upper bound in $N$. Hence, let $\bar{b} \in T_{\lambda} \cap N$ be such that $Y \subseteq \{b \in T_{\lambda} : b \lhd \bar{b}\}$.

There are three cases: either $x = \pi_\alpha(\bar{b})$ or $x \lhd \pi_\alpha(\bar{b})$ or $\pi_\alpha(\bar{b}) \lhd x$. We only consider the latter case, as the other two are analogous.  By (Q4), there exists $\gamma < \alpha$ such that 
\[
\pi_\gamma(\bar{b}) \lhd \pi_\gamma(x) \text{ and } \pi_\alpha(\bar{b}) = Q(\pi_\gamma(\bar{b}),  x).
\]
Let $\delta = \max(\beta, \gamma) + 1$ and
\begin{align*}
I &\coloneqq \big\{z \in T_\delta : z \lhd \pi_\delta (\bar{b}) \text{ and } Q(z, \bar{b}) \in Y\big\},\\X &\coloneqq I \cap D(\pi_\delta(x), \beta, \xi).
\end{align*}
Since $X$ belongs to $N$ by elementarity (using clause (4) of $(\dagger)$ if $\xi = 0$), we can fix $0 < \zeta < \alpha$ such that $X^\delta_\zeta = X$. Both (1) and (2) quickly follow once we show $K = D(\bar{b}, \delta, \zeta)$.

First pick $b \in D(\bar{b}, \delta, \zeta)$, towards showing that $b \in Y \cap N$ and $\pi_\alpha(b) \in D(x, \beta, \xi)$. By the choice of $\zeta$, we know that
\begin{enumerate}[label=(\alph*)]
\item $b = Q(\pi_\delta(b), \bar{b})$,
\item $\pi_\delta(b) \in I$,
\item $\pi_\beta(b) \in X^\beta_\xi$,
\item $\pi_\beta(b) \trianglelefteq \pi_\beta(x)$,
\item $\pi_\delta(b) = Q(\pi_\beta(b), \pi_\delta(x))$.
\end{enumerate}
It follows directly from (a) and (b) that $b \in Y$. Moreover, as $\pi_\delta(b) \in N$---indeed, $T \upharpoonright \alpha \subseteq N$---we also have $b \in N$ by (a) and the elementarity of $N$. Hence, $b \in Y \cap N$, and we are left to show that $\pi_\alpha(b) \in D(x, \beta, \xi)$. The following holds:
\begin{align*}
\pi_\alpha(b) &= \pi_\alpha(Q(\pi_\delta(b), \bar{b}))\\&= Q(\pi_\delta(b), \pi_\alpha(\bar{b}))\\&= Q(\pi_\delta(b), Q(\pi_\delta(\bar{b}), x)) \\&= Q(\pi_\delta(b), x) \\&= Q(Q(\pi_\beta(b), \pi_\delta(x)), x) \\&= Q(\pi_\beta(b), x) \in D(x, \beta, \xi).
\end{align*} 
The first equality follows from (a); the second one from Lemma~\ref{lemma:basicquag}(3); the third one from the choice of $\gamma$ and from Lemma~\ref{lemma:basicquag}(4); the fourth one from (Q3) and from $\pi_\delta(b) \lhd \pi_\delta(\bar{b})$, which holds by (b); the fifth one from (e); the last one from (Q2), and the final membership relation holds by (c) and (d).

So we have shown $D(\bar{b}, \delta, \zeta) \subseteq K$. Now let us pick $b \in Y \cap N$ such that $\pi_\alpha(b) \in D(x, \beta, \xi)$ and let us prove $b \in D(\bar{b}, \delta, \zeta)$. From our assumption we get that 
\begin{enumerate}[label=(\roman*)]
\item $\pi_\beta(b) \trianglelefteq \pi_\beta(x)$,
\item $\pi_\beta(b) \in X^\beta_\xi$,
\item $\pi_\alpha(b) = Q(\pi_\beta(b), x)$.
\end{enumerate}
Moreover, the following holds:
\begin{subclaim}\label{claim:5-1}
$\pi_\delta(b) \lhd \pi_\delta(\bar{b})$.
\end{subclaim}
\begin{proof}
Note first that $\pi_\alpha(b) \lhd \pi_\alpha(\bar{b})$. Indeed, using (Q4), the elementarity of $N$, and the fact that both $b$ and $\bar{b}$ belong to $N$, there exists $\eta < \alpha$ such that $\pi_\eta(b) \lhd \pi_\eta(\bar{b})$ and $b = Q(\pi_\eta(b), \bar{b})$. Then, $\pi_\alpha(b) \lhd \pi_\alpha(\bar{b})$ follows from Lemma~\ref{lemma:basicquag}(4). 

Now, by (iii), by the choice of $\gamma$, and by Lemma~\ref{lemma:basicquag}(4), we have that $\pi_\alpha(b) = Q(\pi_\delta(b), x)$ and $\pi_\alpha(\bar{b}) = Q(\pi_\delta(\bar{b}), x)$. In particular, since the map $Q(\cdot, x) \upharpoonright T_\delta$ is an order-embedding  (Lemma~\ref{lemma:basicquag}(2)), we get that $\pi_\delta(b) \lhd \pi_\delta(\bar{b})$, as we wanted to show.
\end{proof}

By (i)-(iii) and Lemma~\ref{lemma:basicquag}(3), we have $\pi_\delta(b) \in D(\pi_\delta(x), \beta, \xi)$. Moreover, by hypothesis, $b \in Y$ and, by Claim~\ref{claim:5-1}, $\pi_\delta(b) \lhd \pi_\delta(\bar{b})$. Hence, once we prove $b = Q(\pi_\delta(b), \bar{b})$, the fact that $b \in Y$ will imply $\pi_\delta(b) \in I$, and therefore $\pi_\delta(b) \in X = X^\delta_\zeta$. Thus, it remains to show $b = Q(\pi_\delta(b), \bar{b})$. 

As $b \in N$ by hypothesis and $\pi_{\alpha \lambda} \upharpoonright N$ is injective by elementarity of $N$, we just need to show that $\pi_\alpha(b) = \pi_\alpha(Q(\pi_\delta(b), \bar{b}))$---note that $Q(\pi_\delta(b), \bar{b})$ belongs to $N$ by elementarity, since both $\bar{b}$ and $\pi_\delta(b)$ belong to $N$. We have
\begin{align*}
\pi_\alpha(Q(\pi_\delta(b), \bar{b})) &= Q(\pi_\delta(b), \pi_\alpha(\bar{b}))\\&= Q(\pi_\delta(b), Q(\pi_\delta(\bar{b}), x)) \\&= Q(\pi_\delta(b), x) \\&= \pi_\alpha(b),
\end{align*}
where the first equality follows from Lemma~\ref{lemma:basicquag}(3); the second one from the choice of $\gamma$ and Lemma~\ref{lemma:basicquag}(4); the third one holds by (Q3) and Claim~\ref{claim:5-1}, and the last one by (iii) and Lemma~\ref{lemma:basicquag}(4). 

Overall, we have shown $K = D(\bar{b}, \delta, \zeta)$. Now, $K \in N$ clearly follows, since $D(\bar{b}, \delta, \zeta)$ belongs to $N$ by elementarity of $N$.  Moreover, it follows that $D(x, \beta, \xi) \cap X^\alpha_0 \in V_\alpha$. Indeed, the following holds:
\begin{align*}
D(x, \beta, \xi) \cap X^\alpha_0 &= D(x, \beta, \xi) \cap \pi_\alpha[Y \cap N]\\&= \pi_\alpha[K]\\&= \pi_\alpha[D(\bar{b}, \delta , \zeta)] \\&= D(\pi_\alpha(\bar{b}), \delta, \zeta) \in V_\alpha,
\end{align*} 
where the first equality holds by definition of $X^\alpha_0$; the second equality directly follows from the definition of $K$; in the third equality we used $K = D(\bar{b}, \delta, \zeta)$, and the last one holds by Lemma~\ref{lemma:basicquag}(3) and the definition of $D$. This completes the proof.
\end{proof}

\end{document}
